\documentclass[11pt]{article}
\usepackage[T1]{fontenc}
\usepackage[utf8]{inputenc}
\usepackage[spanish]{babel}
\usepackage{amsmath,amssymb}
\usepackage[a4paper,margin=2.3cm]{geometry}
\usepackage{graphicx}
\usepackage[inkscapelatex=false]{svg}
\usepackage{tikz}
\usetikzlibrary{arrows.meta,babel}
\providecommand{\graficanumerica}[2][]{%
  \includesvg[#1]{figuras/#2}%
}
\spanishdecimal{.}
\setlength{\parindent}{0pt}
\setlength{\parskip}{0.35em}
\setlength{\abovedisplayskip}{4pt}
\setlength{\belowdisplayskip}{4pt}
\setlength{\abovedisplayshortskip}{4pt}
\setlength{\belowdisplayshortskip}{4pt}

\newcommand{\encabezado}[3]{%
  {\Large\bfseries #1\par}
  \vspace{2pt}
  {\normalsize #2\hfill #3\par}
  \vspace{6pt}\hrule\vspace{8pt}}

\newcommand{\problema}[2]{%
  \vspace{0.55em}
  \noindent\textbf{#1.}\enspace{\small #2}\par
  \vspace{0.3em}}

\begin{document}



\encabezado{Práctica de repaso: imágenes y polarización de conductores}
{Electrodinámica I}{P.S., J.A., J.R.}

\section*{Ejercicio 1: Carga frente a un Plano Conductor}

Considere una carga puntual $+q$ ubicada sobre el eje $z$ a una distancia $d$ de un plano conductor infinito y conectado a tierra ($V=0$) que ocupa el plano coordenado $z = 0$. El espacio físico de interés es el semiespacio $z > 0$.

\paragraph{Se Pide:}
\begin{enumerate}
    \item \textbf{Método de Imágenes y Potencial:} Utilice el método de las imágenes para plantear la carga imagen necesaria. Escriba la expresión analítica exacta del potencial eléctrico $V(x, y, z)$ en cualquier punto del semiespacio $z > 0$.
    \item \textbf{Fuerza Electrostática:} Calcule la fuerza electrostática neta $\vec{F}$ que experimenta la carga real $+q$ debido a la presencia del plano conductor.
    \item \textbf{Trabajo Externo:} Calcule el trabajo mecánico $W$ que debe realizar un agente externo para mover la carga real a velocidad constante desde su posición original ($z = d$) hasta el doble de esa distancia ($z = 2d$).
    \item \textbf{Densidad Superficial de Carga:} A partir del potencial o del campo eléctrico, determine la densidad superficial de carga inducida $\sigma(x, y)$ sobre el plano conductor ($z=0$). Integre esta densidad sobre todo el plano para comprobar que la carga total inducida es exactamente igual a $-q$.
    \item \textbf{Expansión Multipolar:} Considere el sistema global formado por la carga real y su carga imagen visto desde grandes distancias ($r \gg d$). Escriba la expansión multipolar del potencial a gran distancia y determine cuál es el término dominante (monopolo, dipolo, etc.), calculando explícitamente el momento multipolar efectivo de este sistema respecto al origen.
\end{enumerate}

\begin{figure}[!htbp]
\centering
\begingroup
\shorthandoff{<>}
\resizebox{\textwidth}{!}{%
\begin{tikzpicture}[
  >=Stealth,line cap=round,line join=round,
  every node/.style={font=\small,text=black},scale=0.92,
  campoelectrico/.style={
    draw=black!85,line width=0.55pt
  },
  carga/.style={circle,draw=black,fill=white,inner sep=1.4pt,font=\scriptsize}
]
  % Ambas figuras usan literalmente las mismas curvas y flechas en z>0.
  \def\campoSuperior{
\draw[campoelectrico] plot[smooth,tension=0.45] coordinates {(0.0000,1.0350) (0.0000,2.2000)};
\draw[campoelectrico,-{Stealth[length=1.6mm,width=1.1mm]}] (0.0000,1.6490)--(0.0000,1.8190);
\draw[campoelectrico] plot[smooth,tension=0.45] coordinates {(0.0120,1.0329) (0.2479,1.6557) (0.3647,1.9385) (0.4792,2.2000)};
\draw[campoelectrico,-{Stealth[length=1.6mm,width=1.1mm]}] (0.2495,1.6588)--(0.3138,1.8162);
\draw[campoelectrico] plot[smooth,tension=0.45] coordinates {(-0.0120,1.0329) (-0.2479,1.6557) (-0.3647,1.9385) (-0.4792,2.2000)};
\draw[campoelectrico,-{Stealth[length=1.6mm,width=1.1mm]}] (-0.2495,1.6588)--(-0.3138,1.8162);
\draw[campoelectrico] plot[smooth,tension=0.45] coordinates {(0.0225,1.0268) (0.3288,1.3806) (0.4763,1.5384) (0.6158,1.6789) (0.7595,1.8151) (0.9076,1.9465) (1.0599,2.0730) (1.2238,2.2000)};
\draw[campoelectrico,-{Stealth[length=1.6mm,width=1.1mm]}] (0.6467,1.7088)--(0.7708,1.8250);
\draw[campoelectrico] plot[smooth,tension=0.45] coordinates {(-0.0225,1.0268) (-0.3288,1.3806) (-0.4763,1.5384) (-0.6158,1.6789) (-0.7595,1.8151) (-0.9076,1.9465) (-1.0599,2.0730) (-1.2238,2.2000)};
\draw[campoelectrico,-{Stealth[length=1.6mm,width=1.1mm]}] (-0.6467,1.7088)--(-0.7708,1.8250);
\draw[campoelectrico] plot[smooth,tension=0.45] coordinates {(0.0303,1.0175) (0.2496,1.1416) (0.4408,1.2421) (0.6198,1.3268) (0.7858,1.3963) (0.9550,1.4577) (1.1271,1.5102) (1.3019,1.5530) (1.4789,1.5856) (1.6576,1.6073) (1.8193,1.6171) (1.9992,1.6166) (2.1608,1.6056) (2.3391,1.5811) (2.4976,1.5477) (2.6534,1.5033) (2.8000,1.4497)};
\draw[campoelectrico,-{Stealth[length=1.6mm,width=1.1mm]}] (1.5730,1.5997)--(1.7423,1.6149);
\draw[campoelectrico] plot[smooth,tension=0.45] coordinates {(-0.0303,1.0175) (-0.2496,1.1416) (-0.4408,1.2421) (-0.6198,1.3268) (-0.7858,1.3963) (-0.9550,1.4577) (-1.1271,1.5102) (-1.3019,1.5530) (-1.4789,1.5856) (-1.6576,1.6073) (-1.8193,1.6171) (-1.9992,1.6166) (-2.1608,1.6056) (-2.3391,1.5811) (-2.4976,1.5477) (-2.6534,1.5033) (-2.8000,1.4497)};
\draw[campoelectrico,-{Stealth[length=1.6mm,width=1.1mm]}] (-1.5730,1.5997)--(-1.7423,1.6149);
\draw[campoelectrico] plot[smooth,tension=0.45] coordinates {(0.0345,1.0061) (0.3009,1.0498) (0.4974,1.0735) (0.6950,1.0858) (0.8750,1.0842) (1.0542,1.0684) (1.2313,1.0364) (1.3871,0.9923) (1.5377,0.9328) (1.6808,0.8570) (1.8135,0.7643) (1.9327,0.6548) (2.0244,0.5439) (2.0998,0.4214) (2.1564,0.2891) (2.1918,0.1497) (2.2010,0.0783) (2.2045,0.0064) (2.2045,0.0000)};
\draw[campoelectrico,-{Stealth[length=1.6mm,width=1.1mm]}] (1.5653,0.9207)--(1.7134,0.8372);
\draw[campoelectrico] plot[smooth,tension=0.45] coordinates {(-0.0345,1.0061) (-0.3009,1.0498) (-0.4974,1.0735) (-0.6950,1.0858) (-0.8750,1.0842) (-1.0542,1.0684) (-1.2313,1.0364) (-1.3871,0.9923) (-1.5377,0.9328) (-1.6808,0.8570) (-1.8135,0.7643) (-1.9327,0.6548) (-2.0244,0.5439) (-2.0998,0.4214) (-2.1564,0.2891) (-2.1918,0.1497) (-2.2010,0.0783) (-2.2045,0.0064) (-2.2045,0.0000)};
\draw[campoelectrico,-{Stealth[length=1.6mm,width=1.1mm]}] (-1.5653,0.9207)--(-1.7134,0.8372);
\draw[campoelectrico] plot[smooth,tension=0.45] coordinates {(0.0345,0.9939) (0.3700,0.9278) (0.4920,0.8967) (0.6124,0.8595) (0.7302,0.8148) (0.8283,0.7698) (0.9229,0.7176) (1.0127,0.6578) (1.0966,0.5898) (1.1730,0.5135) (1.2401,0.4290) (1.2876,0.3526) (1.3263,0.2714) (1.3551,0.1862) (1.3732,0.0981) (1.3800,0.0084) (1.3800,0.0000)};
\draw[campoelectrico,-{Stealth[length=1.6mm,width=1.1mm]}] (0.9868,0.6787)--(1.1190,0.5718);
\draw[campoelectrico] plot[smooth,tension=0.45] coordinates {(-0.0345,0.9939) (-0.3700,0.9278) (-0.4920,0.8967) (-0.6124,0.8595) (-0.7302,0.8148) (-0.8283,0.7698) (-0.9229,0.7176) (-1.0127,0.6578) (-1.0966,0.5898) (-1.1730,0.5135) (-1.2401,0.4290) (-1.2876,0.3526) (-1.3263,0.2714) (-1.3551,0.1862) (-1.3732,0.0981) (-1.3800,0.0084) (-1.3800,0.0000)};
\draw[campoelectrico,-{Stealth[length=1.6mm,width=1.1mm]}] (-0.9868,0.6787)--(-1.1190,0.5718);
\draw[campoelectrico] plot[smooth,tension=0.45] coordinates {(0.0303,0.9825) (0.2624,0.8445) (0.4267,0.7341) (0.5532,0.6330) (0.6680,0.5188) (0.7245,0.4488) (0.7650,0.3893) (0.8005,0.3267) (0.8304,0.2613) (0.8540,0.1933) (0.8707,0.1233) (0.8819,0.0160) (0.8821,0.0000)};
\draw[campoelectrico,-{Stealth[length=1.6mm,width=1.1mm]}] (0.6308,0.5627)--(0.7397,0.4322);
\draw[campoelectrico] plot[smooth,tension=0.45] coordinates {(-0.0303,0.9825) (-0.2624,0.8445) (-0.4267,0.7341) (-0.5532,0.6330) (-0.6680,0.5188) (-0.7245,0.4488) (-0.7650,0.3893) (-0.8005,0.3267) (-0.8304,0.2613) (-0.8540,0.1933) (-0.8707,0.1233) (-0.8819,0.0160) (-0.8821,0.0000)};
\draw[campoelectrico,-{Stealth[length=1.6mm,width=1.1mm]}] (-0.6308,0.5627)--(-0.7397,0.4322);
\draw[campoelectrico] plot[smooth,tension=0.45] coordinates {(0.0225,0.9732) (0.1715,0.7927) (0.2702,0.6644) (0.3516,0.5456) (0.4148,0.4366) (0.4675,0.3222) (0.5016,0.2198) (0.5235,0.1141) (0.5316,0.0065) (0.5315,0.0000)};
\draw[campoelectrico,-{Stealth[length=1.6mm,width=1.1mm]}] (0.3789,0.5055)--(0.4569,0.3544);
\draw[campoelectrico] plot[smooth,tension=0.45] coordinates {(-0.0225,0.9732) (-0.1715,0.7927) (-0.2702,0.6644) (-0.3516,0.5456) (-0.4148,0.4366) (-0.4675,0.3222) (-0.5016,0.2198) (-0.5235,0.1141) (-0.5316,0.0065) (-0.5315,0.0000)};
\draw[campoelectrico,-{Stealth[length=1.6mm,width=1.1mm]}] (-0.3789,0.5055)--(-0.4569,0.3544);
\draw[campoelectrico] plot[smooth,tension=0.45] coordinates {(0.0120,0.9671) (0.1315,0.6276) (0.1974,0.4031) (0.2210,0.2977) (0.2386,0.1912) (0.2477,0.1016) (0.2513,0.0117) (0.2513,0.0000)};
\draw[campoelectrico,-{Stealth[length=1.6mm,width=1.1mm]}] (0.1786,0.4787)--(0.2193,0.3137);
\draw[campoelectrico] plot[smooth,tension=0.45] coordinates {(-0.0120,0.9671) (-0.1315,0.6276) (-0.1974,0.4031) (-0.2210,0.2977) (-0.2386,0.1912) (-0.2477,0.1016) (-0.2513,0.0117) (-0.2513,0.0000)};
\draw[campoelectrico,-{Stealth[length=1.6mm,width=1.1mm]}] (-0.1786,0.4787)--(-0.2193,0.3137);
\draw[campoelectrico] plot[smooth,tension=0.45] coordinates {(0.0000,0.9650) (0.0000,0.0110) (0.0000,0.0000)};
\draw[campoelectrico,-{Stealth[length=1.6mm,width=1.1mm]}] (0.0000,0.4710)--(0.0000,0.3010);
  }
  \def\campoInferior{
\draw[campoelectrico] plot[smooth,tension=0.45] coordinates {(0.0000,-2.2000) (0.0000,-1.0350)};
\draw[campoelectrico,-{Stealth[length=1.6mm,width=1.1mm]}] (0.0000,-1.5860)--(0.0000,-1.4160);
\draw[campoelectrico] plot[smooth,tension=0.45] coordinates {(0.4792,-2.2000) (0.3647,-1.9385) (0.2479,-1.6557) (0.0120,-1.0329)};
\draw[campoelectrico,-{Stealth[length=1.6mm,width=1.1mm]}] (0.2213,-1.5824)--(0.1590,-1.4242);
\draw[campoelectrico] plot[smooth,tension=0.45] coordinates {(-0.4792,-2.2000) (-0.3647,-1.9385) (-0.2479,-1.6557) (-0.0120,-1.0329)};
\draw[campoelectrico,-{Stealth[length=1.6mm,width=1.1mm]}] (-0.2213,-1.5824)--(-0.1590,-1.4242);
\draw[campoelectrico] plot[smooth,tension=0.45] coordinates {(1.2238,-2.2000) (1.0599,-2.0730) (0.9076,-1.9465) (0.7595,-1.8151) (0.6158,-1.6789) (0.4763,-1.5384) (0.3288,-1.3806) (0.0225,-1.0268)};
\draw[campoelectrico,-{Stealth[length=1.6mm,width=1.1mm]}] (0.5279,-1.5919)--(0.4101,-1.4693);
\draw[campoelectrico] plot[smooth,tension=0.45] coordinates {(-1.2238,-2.2000) (-1.0599,-2.0730) (-0.9076,-1.9465) (-0.7595,-1.8151) (-0.6158,-1.6789) (-0.4763,-1.5384) (-0.3288,-1.3806) (-0.0225,-1.0268)};
\draw[campoelectrico,-{Stealth[length=1.6mm,width=1.1mm]}] (-0.5279,-1.5919)--(-0.4101,-1.4693);
\draw[campoelectrico] plot[smooth,tension=0.45] coordinates {(2.8000,-1.4497) (2.6534,-1.5033) (2.4976,-1.5477) (2.3391,-1.5811) (2.1608,-1.6056) (1.9992,-1.6166) (1.8193,-1.6171) (1.6576,-1.6073) (1.4789,-1.5856) (1.3019,-1.5530) (1.1271,-1.5102) (0.9550,-1.4577) (0.7858,-1.3963) (0.6198,-1.3268) (0.4408,-1.2421) (0.2496,-1.1416) (0.0303,-1.0175)};
\draw[campoelectrico,-{Stealth[length=1.6mm,width=1.1mm]}] (1.1686,-1.5215)--(1.0053,-1.4744);
\draw[campoelectrico] plot[smooth,tension=0.45] coordinates {(-2.8000,-1.4497) (-2.6534,-1.5033) (-2.4976,-1.5477) (-2.3391,-1.5811) (-2.1608,-1.6056) (-1.9992,-1.6166) (-1.8193,-1.6171) (-1.6576,-1.6073) (-1.4789,-1.5856) (-1.3019,-1.5530) (-1.1271,-1.5102) (-0.9550,-1.4577) (-0.7858,-1.3963) (-0.6198,-1.3268) (-0.4408,-1.2421) (-0.2496,-1.1416) (-0.0303,-1.0175)};
\draw[campoelectrico,-{Stealth[length=1.6mm,width=1.1mm]}] (-1.1686,-1.5215)--(-1.0053,-1.4744);
\draw[campoelectrico] plot[smooth,tension=0.45] coordinates {(2.2045,-0.0000) (2.2045,-0.0064) (2.2010,-0.0783) (2.1918,-0.1497) (2.1564,-0.2891) (2.0998,-0.4214) (2.0244,-0.5439) (1.9327,-0.6548) (1.8135,-0.7643) (1.6808,-0.8570) (1.5377,-0.9328) (1.3871,-0.9923) (1.2313,-1.0364) (1.0542,-1.0684) (0.8750,-1.0842) (0.6950,-1.0858) (0.4974,-1.0735) (0.3009,-1.0498) (0.0345,-1.0061)};
\draw[campoelectrico,-{Stealth[length=1.6mm,width=1.1mm]}] (1.2075,-1.0414)--(1.0400,-1.0702);
\draw[campoelectrico] plot[smooth,tension=0.45] coordinates {(-2.2045,-0.0000) (-2.2045,-0.0064) (-2.2010,-0.0783) (-2.1918,-0.1497) (-2.1564,-0.2891) (-2.0998,-0.4214) (-2.0244,-0.5439) (-1.9327,-0.6548) (-1.8135,-0.7643) (-1.6808,-0.8570) (-1.5377,-0.9328) (-1.3871,-0.9923) (-1.2313,-1.0364) (-1.0542,-1.0684) (-0.8750,-1.0842) (-0.6950,-1.0858) (-0.4974,-1.0735) (-0.3009,-1.0498) (-0.0345,-1.0061)};
\draw[campoelectrico,-{Stealth[length=1.6mm,width=1.1mm]}] (-1.2075,-1.0414)--(-1.0400,-1.0702);
\draw[campoelectrico] plot[smooth,tension=0.45] coordinates {(1.3800,-0.0000) (1.3800,-0.0084) (1.3732,-0.0981) (1.3551,-0.1862) (1.3263,-0.2714) (1.2876,-0.3526) (1.2401,-0.4290) (1.1730,-0.5135) (1.0966,-0.5898) (1.0127,-0.6578) (0.9229,-0.7176) (0.8283,-0.7698) (0.7302,-0.8148) (0.6124,-0.8595) (0.4920,-0.8967) (0.3700,-0.9278) (0.0345,-0.9939)};
\draw[campoelectrico,-{Stealth[length=1.6mm,width=1.1mm]}] (0.8179,-0.7771)--(0.6615,-0.8439);
\draw[campoelectrico] plot[smooth,tension=0.45] coordinates {(-1.3800,-0.0000) (-1.3800,-0.0084) (-1.3732,-0.0981) (-1.3551,-0.1862) (-1.3263,-0.2714) (-1.2876,-0.3526) (-1.2401,-0.4290) (-1.1730,-0.5135) (-1.0966,-0.5898) (-1.0127,-0.6578) (-0.9229,-0.7176) (-0.8283,-0.7698) (-0.7302,-0.8148) (-0.6124,-0.8595) (-0.4920,-0.8967) (-0.3700,-0.9278) (-0.0345,-0.9939)};
\draw[campoelectrico,-{Stealth[length=1.6mm,width=1.1mm]}] (-0.8179,-0.7771)--(-0.6615,-0.8439);
\draw[campoelectrico] plot[smooth,tension=0.45] coordinates {(0.8821,-0.0000) (0.8819,-0.0160) (0.8707,-0.1233) (0.8540,-0.1933) (0.8304,-0.2613) (0.8005,-0.3267) (0.7650,-0.3893) (0.7245,-0.4488) (0.6680,-0.5188) (0.5532,-0.6330) (0.4267,-0.7341) (0.2624,-0.8445) (0.0303,-0.9825)};
\draw[campoelectrico,-{Stealth[length=1.6mm,width=1.1mm]}] (0.5542,-0.6322)--(0.4214,-0.7384);
\draw[campoelectrico] plot[smooth,tension=0.45] coordinates {(-0.8821,-0.0000) (-0.8819,-0.0160) (-0.8707,-0.1233) (-0.8540,-0.1933) (-0.8304,-0.2613) (-0.8005,-0.3267) (-0.7650,-0.3893) (-0.7245,-0.4488) (-0.6680,-0.5188) (-0.5532,-0.6330) (-0.4267,-0.7341) (-0.2624,-0.8445) (-0.0303,-0.9825)};
\draw[campoelectrico,-{Stealth[length=1.6mm,width=1.1mm]}] (-0.5542,-0.6322)--(-0.4214,-0.7384);
\draw[campoelectrico] plot[smooth,tension=0.45] coordinates {(0.5315,-0.0000) (0.5316,-0.0065) (0.5235,-0.1141) (0.5016,-0.2198) (0.4675,-0.3222) (0.4148,-0.4366) (0.3516,-0.5456) (0.2702,-0.6644) (0.1715,-0.7927) (0.0225,-0.9732)};
\draw[campoelectrico,-{Stealth[length=1.6mm,width=1.1mm]}] (0.3487,-0.5510)--(0.2516,-0.6905);
\draw[campoelectrico] plot[smooth,tension=0.45] coordinates {(-0.5315,-0.0000) (-0.5316,-0.0065) (-0.5235,-0.1141) (-0.5016,-0.2198) (-0.4675,-0.3222) (-0.4148,-0.4366) (-0.3516,-0.5456) (-0.2702,-0.6644) (-0.1715,-0.7927) (-0.0225,-0.9732)};
\draw[campoelectrico,-{Stealth[length=1.6mm,width=1.1mm]}] (-0.3487,-0.5510)--(-0.2516,-0.6905);
\draw[campoelectrico] plot[smooth,tension=0.45] coordinates {(0.2513,-0.0000) (0.2513,-0.0117) (0.2477,-0.1016) (0.2386,-0.1912) (0.2210,-0.2977) (0.1974,-0.4031) (0.1315,-0.6276) (0.0120,-0.9671)};
\draw[campoelectrico,-{Stealth[length=1.6mm,width=1.1mm]}] (0.1686,-0.5075)--(0.1174,-0.6696);
\draw[campoelectrico] plot[smooth,tension=0.45] coordinates {(-0.2513,-0.0000) (-0.2513,-0.0117) (-0.2477,-0.1016) (-0.2386,-0.1912) (-0.2210,-0.2977) (-0.1974,-0.4031) (-0.1315,-0.6276) (-0.0120,-0.9671)};
\draw[campoelectrico,-{Stealth[length=1.6mm,width=1.1mm]}] (-0.1686,-0.5075)--(-0.1174,-0.6696);
\draw[campoelectrico] plot[smooth,tension=0.45] coordinates {(0.0000,-0.0000) (0.0000,-0.0110) (0.0000,-0.9650)};
\draw[campoelectrico,-{Stealth[length=1.6mm,width=1.1mm]}] (0.0000,-0.4940)--(0.0000,-0.6640);
  }
  \begin{scope}
    \node[font=\normalsize\bfseries] at (0,2.70)
      {(a) Carga frente al plano};
    \fill[black!3] (-2.80,0) rectangle (2.80,2.20);
    \fill[black!13] (-2.80,-2.20) rectangle (2.80,0);
    \begin{scope}
      \clip (-2.80,0) rectangle (2.80,2.20);
      \campoSuperior
    \end{scope}
    \draw[very thick] (-2.80,0)--(3.35,0);
    \foreach \xx in {-2.35,-1.55,-0.78,0.78,1.55,2.35}{
      \draw[semithick] (\xx,-0.12)++(-0.075,0)--++(0.15,0);
    }
    \node[carga] at (0,1) {$+q$};
    \node[fill=white,inner sep=1pt] at (-2.30,1.97) {$z>0$};
    \node[anchor=west] at (3.50,0) {$z=0$};
    \node[align=center] at (0,-1.08)
      {conductor conectado a tierra\\[3pt]
       $V=0,\quad\vec E=\vec0$};
    \draw[<->] (3.17,0.04)--(3.17,0.98) node[midway,right] {$d$};
    \draw[->,black!55] (-3.70,1.55)--(-3.15,1.55) node[right] {$x$};
    \draw[->,black!55] (-3.70,1.55)--(-3.70,2.10) node[above] {$z$};
  \end{scope}
  \begin{scope}[xshift=7.65cm]
    \node[font=\normalsize\bfseries] at (0,2.70)
      {(b) Dipolo carga--imagen};
    \fill[black!3] (-2.80,0) rectangle (2.80,2.20);
    \begin{scope}
      \clip (-2.80,-2.20) rectangle (2.80,2.20);
      \campoSuperior
      \campoInferior
    \end{scope}
    \draw[densely dashed,black!55] (-2.80,0)--(3.35,0);
    \node[carga] at (0,1) {$+q$};
    \node[carga] at (0,-1) {$-q$};
    \node[fill=white,inner sep=1pt] at (-2.30,1.97) {$z>0$};
    \node[anchor=west] at (3.50,0) {$z=0$};
    \draw[<->] (3.17,0.04)--(3.17,0.98) node[midway,right] {$d$};
    \draw[<->] (3.17,-0.04)--(3.17,-0.98) node[midway,right] {$d$};
    \draw[->,black!55] (-3.70,1.55)--(-3.15,1.55) node[right] {$x$};
    \draw[->,black!55] (-3.70,1.55)--(-3.70,2.10) node[above] {$z$};
  \end{scope}
  \node[align=center] at (3.82,-2.95)
    {En $z>0$ coinciden las líneas de campo y su dirección:\\[3pt]
     $V_{\mathrm{plano}}=V_{\mathrm{imagen}},\qquad
      \vec E_{\mathrm{plano}}=\vec E_{\mathrm{imagen}}$};
\end{tikzpicture}
}
\endgroup
\caption{Líneas de campo en dos cortes del plano $xz$. (a) Salen de
$+q$ y llegan perpendicularmente al conductor; dentro de él el campo
es cero. (b) Salen de $+q$ y terminan en la carga imagen $-q$.
En el semiespacio $z>0$ el trazado es idéntico en ambos sistemas:
el par carga--imagen reproduce exactamente el campo de la carga
frente al plano. Las flechas indican la dirección del campo.}
\label{fig:plano-dipolo}
\end{figure}

\vspace{0.15cm}
\hrule
\vspace{0.15cm}

\subsection*{Solución Detallada}

\subsubsection*{1. Método de Imágenes y Potencial}
La condición de frontera de Dirichlet es
\[
 V(x,y,0)=0.
\]
Para satisfacerla, introducimos una carga imagen $-q$ situada
especularmente en $z=-d$. Las posiciones de las dos fuentes son
\begin{equation}
 \vec{x}\,'_{\mathrm{real}}=d\hat{z},\qquad
 \vec{x}\,'_{\mathrm{ima}}=-d\hat{z}.
 \label{eq:fuentes}
\end{equation}
El punto de observación es $\vec{x}=x\hat{x}+y\hat{y}+z\hat{z}$.
Las primas identifican las posiciones de las fuentes; $\vec{x}$ recorre
el semiespacio físico $z>0$, como en la figura~\ref{fig:plano-dipolo}.

El potencial de una carga puntual situada en $\vec{x}\,'$ es
\begin{equation}
 V(\vec{x})=\frac{1}{4\pi\varepsilon_0}
 \frac{q}{|\vec{x}-\vec{x}\,'|}.
 \label{eq:potencial-puntual}
\end{equation}
Aplicando \eqref{eq:potencial-puntual} a ambas fuentes y sumando sus
contribuciones, proponemos
\begin{equation}
 V(\vec{x})=\frac{1}{4\pi\varepsilon_0}
 \left[\frac{q}{|\vec{x}-\vec{x}\,'_{\mathrm{real}}|}
       -\frac{q}{|\vec{x}-\vec{x}\,'_{\mathrm{ima}}|}\right].
 \label{eq:superposicion-plano}
\end{equation}
Calculemos las separaciones antes de escribir sus módulos:
\begin{align*}
 \vec{x}-\vec{x}\,'_{\mathrm{real}}
 &=x\hat{x}+y\hat{y}+(z-d)\hat{z},\\
 \vec{x}-\vec{x}\,'_{\mathrm{ima}}
 &=x\hat{x}+y\hat{y}+(z+d)\hat{z}.
\end{align*}
Así, las distancias a la carga real y a la imagen son, respectivamente,
\begin{align}
 |\vec{x}-\vec{x}\,'_{\mathrm{real}}|
 &=\sqrt{x^2+y^2+(z-d)^2},\label{eq:distancia-real}\\
 |\vec{x}-\vec{x}\,'_{\mathrm{ima}}|
 &=\sqrt{x^2+y^2+(z+d)^2}.\label{eq:distancia-imagen}
\end{align}
Al sustituir \eqref{eq:distancia-real} y \eqref{eq:distancia-imagen}
en \eqref{eq:superposicion-plano}, obtenemos el potencial exacto:
\begin{equation}
 \boxed{V(x,y,z)=\frac{q}{4\pi\varepsilon_0}
 \left[\frac{1}{\sqrt{x^2+y^2+(z-d)^2}}
       -\frac{1}{\sqrt{x^2+y^2+(z+d)^2}}\right],\quad z>0.}
 \label{eq:plano-potencial}
\end{equation}
La expresión es singular en la posición de la carga real. En el plano,
las dos distancias coinciden y comprobamos explícitamente que
\[
 V(x,y,0)=\frac{q}{4\pi\varepsilon_0}
 \left[\frac{1}{\sqrt{x^2+y^2+d^2}}
       -\frac{1}{\sqrt{x^2+y^2+d^2}}\right]=0.
\]
\textit{La imagen es una fuente auxiliar:} reproduce la condición de
borde y el potencial en $z>0$; la carga física del conductor está
distribuida sobre su superficie. La unicidad del problema con esta
frontera y $V\to0$ en el infinito justifica la construcción.
Por ello, el sistema carga--plano y el dipolo formado por la carga real
y la imagen son equivalentes en el semiespacio positivo, como muestra
la figura~\ref{fig:plano-dipolo}. Esta equivalencia es exacta;
la aproximación por un dipolo puntual se usará solo cuando $r\gg d$.

\subsubsection*{2. Fuerza Electrostática}
La fuerza sobre $+q$ se obtiene del campo producido por el conductor,
representado en $z>0$ por la imagen. Evaluamos ese campo en la posición
de la carga real; su propio campo no participa en esta fuerza.
\begin{equation}
 \vec{E}_{\mathrm{ima}}(\vec{x}\,'_{\mathrm{real}})
 =\frac{1}{4\pi\varepsilon_0}
 \frac{-q}{|\vec{x}\,'_{\mathrm{real}}-\vec{x}\,'_{\mathrm{ima}}|^3}
 (\vec{x}\,'_{\mathrm{real}}-\vec{x}\,'_{\mathrm{ima}}).
 \label{eq:campo-imagen}
\end{equation}
Usando las posiciones de \eqref{eq:fuentes}, el vector de separación
apunta desde la imagen hacia la carga real:
\[
 \vec{x}\,'_{\mathrm{real}}-\vec{x}\,'_{\mathrm{ima}}
 =d\hat{z}-(-d\hat{z})=2d\hat{z},\qquad
 |\vec{x}\,'_{\mathrm{real}}-\vec{x}\,'_{\mathrm{ima}}|=2d.
\]
Sustituyendo esta separación en \eqref{eq:campo-imagen},
\begin{align*}
 \vec{E}_{\mathrm{ima}}(0,0,d)
 &=\frac{1}{4\pi\varepsilon_0}\frac{-q}{(2d)^3}(2d\hat{z})\\
 &=-\frac{q}{4\pi\varepsilon_0(2d)^2}\hat{z}
 =-\frac{q}{16\pi\varepsilon_0d^2}\hat{z}.
\end{align*}
Multiplicamos por la carga real:
\begin{equation}
 \boxed{\vec{F}_{\mathrm{elec}}(d)
 =q\vec{E}_{\mathrm{ima}}(0,0,d)
 =-\frac{q^2}{16\pi\varepsilon_0d^2}\hat{z}.}
 \label{eq:plano-fuerza}
\end{equation}
El signo negativo indica una fuerza atractiva, dirigida hacia el plano.

\subsubsection*{3. Trabajo Externo}
Al mover la carga hasta una altura $z$, su imagen pasa a $-z$: la
separación entre ambas es $2z$. En consecuencia,
\begin{equation}
 \vec{F}_{\mathrm{elec}}(z)
 =-\frac{q^2}{16\pi\varepsilon_0z^2}\hat{z}.
 \label{eq:fuerza-altura}
\end{equation}
Para que la velocidad sea constante, el agente externo debe compensar
la fuerza de \eqref{eq:fuerza-altura} en cada posición:
\[
 \vec{F}_{\mathrm{ext}}=-\vec{F}_{\mathrm{elec}}
 =\frac{q^2}{16\pi\varepsilon_0z^2}\hat{z}.
\]
El desplazamiento es $d\vec{l}=dz\,\hat{z}$. Escribimos el producto
escalar antes de integrar:
\begin{align*}
 W&=\int_d^{2d}\vec{F}_{\mathrm{ext}}\cdot d\vec{l}\\
 &=\int_d^{2d}\frac{q^2}{16\pi\varepsilon_0z^2}
    \underbrace{\hat{z}\cdot\hat{z}}_{1}\,dz
 =\frac{q^2}{16\pi\varepsilon_0}\int_d^{2d}z^{-2}\,dz\\
 &=\frac{q^2}{16\pi\varepsilon_0}
    \left[-\frac1z\right]_d^{2d}\\
 &=\frac{q^2}{16\pi\varepsilon_0}
    \left(-\frac{1}{2d}+\frac1d\right)
 =\frac{q^2}{16\pi\varepsilon_0}\frac{-1+2}{2d}.
\end{align*}
Así,
\begin{equation}
 \boxed{W=\frac{q^2}{32\pi\varepsilon_0d}>0.}
 \label{eq:plano-trabajo}
\end{equation}
El trabajo es positivo porque el agente externo aleja la carga en
contra de la atracción del conductor.

\subsubsection*{4. Densidad Superficial de Carga y Carga Total Inducida}
\textit{Condición de salto normal.} Elegimos la normal
$\hat{n}=\hat{z}$, dirigida desde el conductor hacia $z>0$. La condición
de frontera del campo es
\begin{equation}
 \hat{n}\cdot(\vec{E}_{\mathrm{ext}}-\vec{E}_{\mathrm{int}})
 =\frac{\sigma}{\varepsilon_0}.
 \label{eq:salto-normal}
\end{equation}
En equilibrio electrostático, $\vec{E}_{\mathrm{int}}=\vec{0}$; la
conexión a tierra fija además el potencial del conductor en cero.
Usando $\vec{E}=-\nabla V$, la condición anterior se reduce a
\[
 \sigma=\varepsilon_0 E_{\perp,\mathrm{ext}}
 =\varepsilon_0 E_z(0^+)
 =-\varepsilon_0\left.\frac{\partial V}{\partial z}\right|_{z=0^+}.
\]
El símbolo $0^+$ indica que evaluamos el campo desde el vacío.

Derivamos \eqref{eq:plano-potencial} con respecto a $z$. Para cada
denominador, la regla de la cadena entrega
\[
 \frac{\partial}{\partial z}
 [x^2+y^2+(z\mp d)^2]^{-1/2}
 =-\frac{z\mp d}{[x^2+y^2+(z\mp d)^2]^{3/2}}.
\]
Por tanto,
\begin{equation}
 \frac{\partial V}{\partial z}
 =\frac{q}{4\pi\varepsilon_0}
 \left[-\frac{z-d}{[x^2+y^2+(z-d)^2]^{3/2}}
       +\frac{z+d}{[x^2+y^2+(z+d)^2]^{3/2}}\right].
 \label{eq:derivada-plano}
\end{equation}
Evaluando \eqref{eq:derivada-plano} en $z=0$ y escribiendo
$s^2=x^2+y^2$, las dos contribuciones se suman:
\[
 \left.\frac{\partial V}{\partial z}\right|_{0^+}
 =\frac{q}{4\pi\varepsilon_0}
 \left[\frac{d}{(s^2+d^2)^{3/2}}
       +\frac{d}{(s^2+d^2)^{3/2}}\right]
 =\frac{2qd}{4\pi\varepsilon_0(s^2+d^2)^{3/2}}.
\]
Al sustituir en la condición de salto \eqref{eq:salto-normal}, obtenemos
\begin{equation}
 \boxed{\sigma(s)=-\frac{qd}{2\pi(s^2+d^2)^{3/2}}.}
 \label{eq:plano-sigma}
\end{equation}
La densidad es negativa y su magnitud es mayor debajo de la carga real.

\textit{Carga total inducida.} $Q_{\mathrm{ind}}$ es la carga física
inducida sobre todo el plano conductor infinito, distribuida según
$\sigma(s)$. Con $dA=s\,ds\,d\phi$, la integral angular da $2\pi s\,ds$:
\begin{align*}
 Q_{\mathrm{ind}}
 &=\int_0^{2\pi}\!\int_0^\infty\sigma(s)s\,ds\,d\phi\\
 &=\int_0^\infty
   \left[-\frac{qd}{2\pi(s^2+d^2)^{3/2}}\right]2\pi s\,ds
 =-qd\int_0^\infty\frac{s\,ds}{(s^2+d^2)^{3/2}}.
\end{align*}
Hacemos el cambio $u=s^2+d^2$, con $du=2s\,ds$.
\textbf{También cambian los límites de integración:}
\[
 s=0\ \Longrightarrow\ u=d^2,\qquad
 s\to\infty\ \Longrightarrow\ u\to\infty.
\]
La integral queda
\begin{align*}
 Q_{\mathrm{ind}}
 &=-\frac{qd}{2}\int_{d^2}^\infty u^{-3/2}\,du\\
 &=-\frac{qd}{2}
   \left[\frac{u^{-1/2}}{-1/2}\right]_{d^2}^{\infty}
 =qd\left[\frac{1}{\sqrt{u}}\right]_{d^2}^{\infty}\\
 &=qd\left(0-\frac{1}{\sqrt{d^2}}\right)
 =qd\left(0-\frac1d\right),\qquad d>0.
\end{align*}
Por lo tanto,
\begin{equation}
 \boxed{Q_{\mathrm{ind}}=-q.}
 \label{eq:plano-carga-total}
\end{equation}

\subsubsection*{5. Expansión Multipolar}
Para $r=|\vec{x}|\gg d$, con $\theta$ medido desde el eje $z$, el
potencial se descompone en contribuciones de distinto orden, siguiendo
la definición del apunte de Electrodinámica de Guillermo Rubilar:
\begin{equation}
 V(\vec{x})=\sum_{n=0}^{\infty}V^{(n)}(\vec{x})
 =V^{(0)}(\vec{x})+V^{(1)}(\vec{x})+\cdots.
 \label{eq:expansion-multipolar}
\end{equation}
La expansión se aplica al par auxiliar de \textbf{dos cargas puntuales},
real e imagen. Sus momentos se calculan sumando sus contribuciones:
\begin{align}
 V^{(0)}(\vec{x})&=\frac{1}{4\pi\varepsilon_0}\frac{Q}{r},
 &Q&=\sum_a q_a,\label{eq:def-monopolo}\\
 V^{(1)}(\vec{x})&=\frac{1}{4\pi\varepsilon_0}
 \frac{\vec{p}\cdot\vec{x}}{r^3},
 &\vec{p}&=\sum_a q_a\vec{x}\,'_a.\label{eq:potencial-dipolar}
\end{align}
La suma recorre ambas fuentes.

\textit{Monopolo.} Sumamos la carga real y la carga física inducida
sobre el plano, calculada en \eqref{eq:plano-carga-total}:
\begin{align*}
 Q&=q+Q_{\mathrm{ind}}\\
  &=q+(-q)=0.
\end{align*}
El par auxiliar tiene la misma carga total, pues
$q_{\mathrm{ima}}=Q_{\mathrm{ind}}=-q$. Así, $V^{(0)}(\vec{x})=0$.

\textit{Dipolo.} Sumamos los momentos de las dos cargas puntuales:
\begin{equation}
 \begin{aligned}
 \vec{p}
 &=q\vec{x}\,'_{\mathrm{real}}+(-q)\vec{x}\,'_{\mathrm{ima}}\\
 &=q(d\hat{z})+(-q)(-d\hat{z})=2qd\hat{z}.
 \end{aligned}
 \label{eq:momento-dipolar}
\end{equation}
Este momento apunta de la carga negativa hacia la positiva.

Truncamos \eqref{eq:expansion-multipolar} en el dipolo dominante.
Para evaluarlo, recordamos
\[
 \hat{r}=\sin\theta\cos\phi\,\hat{x}
          +\sin\theta\sin\phi\,\hat{y}
          +\cos\theta\,\hat{z}.
\]
Como $\vec{x}=r\hat{r}$ y los vectores unitarios son ortogonales,
\[
 \hat{z}\cdot\hat{r}=\cos\theta,\qquad
 \vec{p}\cdot\vec{x}=2qd\,r\cos\theta.
\]
Al sustituir en \eqref{eq:potencial-dipolar}, queda
\begin{equation}
 \boxed{V(r,\theta)\simeq
 \frac{2qd\cos\theta}{4\pi\varepsilon_0r^2},\qquad r\gg d.}
 \label{eq:plano-dipolo}
\end{equation}
Los términos restantes son menores para $d/r\ll1$ y decrecen más rápido
que $1/r^2$. A grandes distancias, el potencial y el campo se ven como
los de un dipolo. La figura~\ref{fig:numerica-plano}(d) muestra cómo
disminuye el error al conservar únicamente ese término dominante.

La visualización numérica del ejercicio 1 permite ver estos resultados:
la figura~\ref{fig:numerica-plano}(a) muestra el potencial nulo en el
plano, (b) las líneas del campo que reproduce el sistema carga--imagen,
y (c) la carga superficial negativa concentrada bajo la carga real.

\newpage

\section*{Ejercicio 2: Cilindro Conductor en un Campo Uniforme}

Un cilindro circular largo y neutro de radio $R$, hecho de material conductor ideal, se coloca en el seno de un campo eléctrico externo uniforme $\vec{E}_0 = E_0 \hat{x}$ (es decir, el eje del cilindro coincide con el eje $z$). Determine la densidad de carga superficial inducida en el cilindro.

\begin{figure}[!htbp]
\centering
\begin{tikzpicture}[
  >=Stealth,line cap=round,line join=round,
  every node/.style={font=\small,text=black},scale=1.45,
  campototal/.style={draw=black!85,line width=0.55pt}
]
  % Curvas del campo total calculadas con el potencial exterior.
  \begin{scope}
    \clip (-2.80,-2.10) rectangle (2.80,2.10);
\draw[campototal] plot[smooth,tension=0.45] coordinates {(-2.8000,0.0000) (-1.0000,0.0000) (-1.0000,0.0000)};
\draw[campototal,-{Stealth[length=1.7mm,width=1.1mm]}] (-1.8300,0.0000)--(-1.6100,0.0000);
\draw[campototal] plot[smooth,tension=0.45] coordinates {(1.0000,0.0000) (1.0000,0.0000) (2.8000,0.0000)};
\draw[campototal,-{Stealth[length=1.7mm,width=1.1mm]}] (1.9700,0.0000)--(2.1900,0.0000);
\draw[campototal] plot[smooth,tension=0.45] coordinates {(-2.8000,0.3200) (-2.2064,0.2998) (-1.7212,0.2710) (-1.3272,0.2323) (-0.9891,0.1811) (-0.9836,0.1801)};
\draw[campototal,-{Stealth[length=1.7mm,width=1.1mm]}] (-1.8169,0.2782)--(-1.5976,0.2611);
\draw[campototal] plot[smooth,tension=0.45] coordinates {(0.9836,0.1801) (0.9891,0.1811) (1.3272,0.2323) (1.7212,0.2710) (2.2064,0.2998) (2.8000,0.3200)};
\draw[campototal,-{Stealth[length=1.7mm,width=1.1mm]}] (1.9614,0.2861)--(2.1811,0.2975);
\draw[campototal] plot[smooth,tension=0.45] coordinates {(-2.8000,-0.3200) (-2.2064,-0.2998) (-1.7212,-0.2710) (-1.3272,-0.2323) (-0.9891,-0.1811) (-0.9836,-0.1801)};
\draw[campototal,-{Stealth[length=1.7mm,width=1.1mm]}] (-1.8169,-0.2782)--(-1.5976,-0.2611);
\draw[campototal] plot[smooth,tension=0.45] coordinates {(0.9836,-0.1801) (0.9891,-0.1811) (1.3272,-0.2323) (1.7212,-0.2710) (2.2064,-0.2998) (2.8000,-0.3200)};
\draw[campototal,-{Stealth[length=1.7mm,width=1.1mm]}] (1.9614,-0.2861)--(2.1811,-0.2975);
\draw[campototal] plot[smooth,tension=0.45] coordinates {(-2.8000,0.6400) (-2.4585,0.6210) (-2.1713,0.5994) (-1.9027,0.5726) (-1.6707,0.5419) (-1.4579,0.5053) (-1.2646,0.4625) (-1.0913,0.4138) (-0.9383,0.3607) (-0.9334,0.3588)};
\draw[campototal,-{Stealth[length=1.7mm,width=1.1mm]}] (-1.7764,0.5578)--(-1.5589,0.5249);
\draw[campototal] plot[smooth,tension=0.45] coordinates {(0.9334,0.3588) (0.9383,0.3607) (1.0913,0.4138) (1.2646,0.4625) (1.4579,0.5053) (1.6707,0.5419) (1.9027,0.5726) (2.1713,0.5994) (2.4585,0.6210) (2.8000,0.6400)};
\draw[campototal,-{Stealth[length=1.7mm,width=1.1mm]}] (1.9345,0.5759)--(2.1534,0.5975);
\draw[campototal] plot[smooth,tension=0.45] coordinates {(-2.8000,-0.6400) (-2.4585,-0.6210) (-2.1713,-0.5994) (-1.9027,-0.5726) (-1.6707,-0.5419) (-1.4579,-0.5053) (-1.2646,-0.4625) (-1.0913,-0.4138) (-0.9383,-0.3607) (-0.9334,-0.3588)};
\draw[campototal,-{Stealth[length=1.7mm,width=1.1mm]}] (-1.7764,-0.5578)--(-1.5589,-0.5249);
\draw[campototal] plot[smooth,tension=0.45] coordinates {(0.9334,-0.3588) (0.9383,-0.3607) (1.0913,-0.4138) (1.2646,-0.4625) (1.4579,-0.5053) (1.6707,-0.5419) (1.9027,-0.5726) (2.1713,-0.5994) (2.4585,-0.6210) (2.8000,-0.6400)};
\draw[campototal,-{Stealth[length=1.7mm,width=1.1mm]}] (1.9345,-0.5759)--(2.1534,-0.5975);
\draw[campototal] plot[smooth,tension=0.45] coordinates {(-2.8000,0.9600) (-2.4410,0.9330) (-2.1187,0.8998) (-1.8335,0.8599) (-1.5859,0.8135) (-1.3589,0.7568) (-1.1709,0.6948) (-1.0055,0.6240) (-0.8481,0.5368) (-0.8450,0.5348)};
\draw[campototal,-{Stealth[length=1.7mm,width=1.1mm]}] (-1.7055,0.8387)--(-1.4904,0.7929);
\draw[campototal] plot[smooth,tension=0.45] coordinates {(0.8450,0.5348) (0.8481,0.5368) (1.0055,0.6240) (1.1709,0.6948) (1.3589,0.7568) (1.5859,0.8135) (1.8335,0.8599) (2.1187,0.8998) (2.4410,0.9330) (2.8000,0.9600)};
\draw[campototal,-{Stealth[length=1.7mm,width=1.1mm]}] (1.8872,0.8678)--(2.1052,0.8975);
\draw[campototal] plot[smooth,tension=0.45] coordinates {(-2.8000,-0.9600) (-2.4410,-0.9330) (-2.1187,-0.8998) (-1.8335,-0.8599) (-1.5859,-0.8135) (-1.3589,-0.7568) (-1.1709,-0.6948) (-1.0055,-0.6240) (-0.8481,-0.5368) (-0.8450,-0.5348)};
\draw[campototal,-{Stealth[length=1.7mm,width=1.1mm]}] (-1.7055,-0.8387)--(-1.4904,-0.7929);
\draw[campototal] plot[smooth,tension=0.45] coordinates {(0.8450,-0.5348) (0.8481,-0.5368) (1.0055,-0.6240) (1.1709,-0.6948) (1.3589,-0.7568) (1.5859,-0.8135) (1.8335,-0.8599) (2.1187,-0.8998) (2.4410,-0.9330) (2.8000,-0.9600)};
\draw[campototal,-{Stealth[length=1.7mm,width=1.1mm]}] (1.8872,-0.8678)--(2.1052,-0.8975);
\draw[campototal] plot[smooth,tension=0.45] coordinates {(-2.8000,1.3000) (-2.3697,1.2614) (-2.0126,1.2164) (-1.6934,1.1610) (-1.4307,1.0988) (-1.1910,1.0215) (-1.0910,0.9807) (-0.9936,0.9341) (-0.8997,0.8808) (-0.8251,0.8304) (-0.7548,0.7742) (-0.7024,0.7249) (-0.6959,0.7182)};
\draw[campototal,-{Stealth[length=1.7mm,width=1.1mm]}] (-1.5879,1.1383)--(-1.3749,1.0832);
\draw[campototal] plot[smooth,tension=0.45] coordinates {(0.6959,0.7182) (0.7024,0.7249) (0.7548,0.7742) (0.8251,0.8304) (0.8997,0.8808) (0.9936,0.9341) (1.0910,0.9807) (1.1910,1.0215) (1.4307,1.0988) (1.6934,1.1610) (2.0126,1.2164) (2.3697,1.2614) (2.8000,1.3000)};
\draw[campototal,-{Stealth[length=1.7mm,width=1.1mm]}] (1.8082,1.1822)--(2.0253,1.2174);
\draw[campototal] plot[smooth,tension=0.45] coordinates {(-2.8000,-1.3000) (-2.3697,-1.2614) (-2.0126,-1.2164) (-1.6934,-1.1610) (-1.4307,-1.0988) (-1.1910,-1.0215) (-1.0910,-0.9807) (-0.9936,-0.9341) (-0.8997,-0.8808) (-0.8251,-0.8304) (-0.7548,-0.7742) (-0.7024,-0.7249) (-0.6959,-0.7182)};
\draw[campototal,-{Stealth[length=1.7mm,width=1.1mm]}] (-1.5879,-1.1383)--(-1.3749,-1.0832);
\draw[campototal] plot[smooth,tension=0.45] coordinates {(0.6959,-0.7182) (0.7024,-0.7249) (0.7548,-0.7742) (0.8251,-0.8304) (0.8997,-0.8808) (0.9936,-0.9341) (1.0910,-0.9807) (1.1910,-1.0215) (1.4307,-1.0988) (1.6934,-1.1610) (2.0126,-1.2164) (2.3697,-1.2614) (2.8000,-1.3000)};
\draw[campototal,-{Stealth[length=1.7mm,width=1.1mm]}] (1.8082,-1.1822)--(2.0253,-1.2174);
\draw[campototal] plot[smooth,tension=0.45] coordinates {(-2.8000,1.6000) (-2.3162,1.5541) (-1.8882,1.4961) (-1.5345,1.4288) (-1.3771,1.3908) (-1.2213,1.3464) (-1.0848,1.3004) (-0.9512,1.2469) (-0.8377,1.1923) (-0.7441,1.1384) (-0.6558,1.0763) (-0.5882,1.0170) (-0.5284,0.9498) (-0.4881,0.8902) (-0.4807,0.8769)};
\draw[campototal,-{Stealth[length=1.7mm,width=1.1mm]}] (-1.4259,1.4040)--(-1.2140,1.3450);
\draw[campototal] plot[smooth,tension=0.45] coordinates {(0.4807,0.8769) (0.4881,0.8902) (0.5284,0.9498) (0.5882,1.0170) (0.6558,1.0763) (0.7441,1.1384) (0.8377,1.1923) (0.9512,1.2469) (1.0848,1.3004) (1.2213,1.3464) (1.3771,1.3908) (1.5345,1.4288) (1.8882,1.4961) (2.3162,1.5541) (2.8000,1.6000)};
\draw[campototal,-{Stealth[length=1.7mm,width=1.1mm]}] (1.6996,1.4620)--(1.9164,1.4996);
\draw[campototal] plot[smooth,tension=0.45] coordinates {(-2.8000,-1.6000) (-2.3162,-1.5541) (-1.8882,-1.4961) (-1.5345,-1.4288) (-1.3771,-1.3908) (-1.2213,-1.3464) (-1.0848,-1.3004) (-0.9512,-1.2469) (-0.8377,-1.1923) (-0.7441,-1.1384) (-0.6558,-1.0763) (-0.5882,-1.0170) (-0.5284,-0.9498) (-0.4881,-0.8902) (-0.4807,-0.8769)};
\draw[campototal,-{Stealth[length=1.7mm,width=1.1mm]}] (-1.4259,-1.4040)--(-1.2140,-1.3450);
\draw[campototal] plot[smooth,tension=0.45] coordinates {(0.4807,-0.8769) (0.4881,-0.8902) (0.5284,-0.9498) (0.5882,-1.0170) (0.6558,-1.0763) (0.7441,-1.1384) (0.8377,-1.1923) (0.9512,-1.2469) (1.0848,-1.3004) (1.2213,-1.3464) (1.3771,-1.3908) (1.5345,-1.4288) (1.8882,-1.4961) (2.3162,-1.5541) (2.8000,-1.6000)};
\draw[campototal,-{Stealth[length=1.7mm,width=1.1mm]}] (1.6996,-1.4620)--(1.9164,-1.4996);
\draw[campototal] plot[smooth,tension=0.45] coordinates {(-2.8000,1.9000) (-2.3698,1.8610) (-1.9944,1.8163) (-1.6564,1.7647) (-1.3383,1.7029) (-1.0585,1.6350) (-0.7998,1.5579) (-0.5960,1.4863) (-0.2613,1.3538) (-0.1593,1.3183) (-0.0718,1.2977) (-0.0001,1.2919) (0.0716,1.2977) (0.1591,1.3183) (0.2611,1.3537) (0.5958,1.4863) (0.7996,1.5579) (1.0583,1.6349) (1.3382,1.7029) (1.6562,1.7647) (1.9943,1.8163) (2.3696,1.8609) (2.8000,1.9000)};
\draw[campototal,-{Stealth[length=1.7mm,width=1.1mm]}] (0.4432,1.4268)--(0.6484,1.5061);
\draw[campototal] plot[smooth,tension=0.45] coordinates {(-2.8000,-1.9000) (-2.3698,-1.8610) (-1.9944,-1.8163) (-1.6564,-1.7647) (-1.3383,-1.7029) (-1.0585,-1.6350) (-0.7998,-1.5579) (-0.5960,-1.4863) (-0.2613,-1.3538) (-0.1593,-1.3183) (-0.0718,-1.2977) (-0.0001,-1.2919) (0.0716,-1.2977) (0.1591,-1.3183) (0.2611,-1.3537) (0.5958,-1.4863) (0.7996,-1.5579) (1.0583,-1.6349) (1.3382,-1.7029) (1.6562,-1.7647) (1.9943,-1.8163) (2.3696,-1.8609) (2.8000,-1.9000)};
\draw[campototal,-{Stealth[length=1.7mm,width=1.1mm]}] (0.4432,-1.4268)--(0.6484,-1.5061);
\draw[campototal] plot[smooth,tension=0.45] coordinates {(-2.8000,2.1000) (-2.1906,2.0437) (-1.6552,1.9740) (-1.3888,1.9301) (-1.1412,1.8830) (-0.3322,1.7070) (-0.1540,1.6818) (0.0077,1.6741) (0.1694,1.6833) (0.3474,1.7097) (1.1387,1.8825) (1.3863,1.9296) (1.6526,1.9736) (2.1881,2.0434) (2.8000,2.1000)};
\draw[campototal,-{Stealth[length=1.7mm,width=1.1mm]}] (0.4527,1.7321)--(0.6673,1.7802);
\draw[campototal] plot[smooth,tension=0.45] coordinates {(-2.8000,-2.1000) (-2.1906,-2.0437) (-1.6552,-1.9740) (-1.3888,-1.9301) (-1.1412,-1.8830) (-0.3322,-1.7070) (-0.1540,-1.6818) (0.0077,-1.6741) (0.1694,-1.6833) (0.3474,-1.7097) (1.1387,-1.8825) (1.3863,-1.9296) (1.6526,-1.9736) (2.1881,-2.0434) (2.8000,-2.1000)};
\draw[campototal,-{Stealth[length=1.7mm,width=1.1mm]}] (0.4527,-1.7321)--(0.6673,-1.7802);
  \end{scope}

  % El campo aplicado es uniforme; se muestra por separado arriba.
  \foreach \xx in {-2.55,-0.65,1.25}{
    \draw[->,thick,black!65,densely dashed] (\xx,2.40)--++(1.35,0);
  }
  \node[above,font=\normalsize] at (0,2.60)
    {$\vec E_0=E_0\hat{x}$: campo externo uniforme};

  % Banda superficial continua: el sombreado sigue |cos(phi)|.
  \fill[black!5] (0,0) circle (1);
  \foreach \ang in {0,4,...,356}{
    \pgfmathsetmacro{\sectorEnd}{\ang+4.15}
    \pgfmathtruncatemacro{\sectorGray}{25*abs(cos(\ang+2))}
    \fill[black!\sectorGray]
      (\ang:1) arc[start angle=\ang,end angle=\sectorEnd,radius=1]
      --(\sectorEnd:0.84)
      arc[start angle=\sectorEnd,end angle=\ang,radius=0.84]--cycle;
  }
  \draw[thick] (0,0) circle (1);
  \foreach \ang in {-60,-30,0,30,60}{
    \node[font=\scriptsize] at (\ang:0.91) {$+$};
  }
  \foreach \ang in {120,150,180,210,240}{
    \node[font=\scriptsize] at (\ang:0.91) {$-$};
  }

  % Estos dos vectores representan las contribuciones, no el campo neto.
  \draw[->,semithick,black!65,densely dashed] (-0.40,0.30)--(0.40,0.30);
  \node at (0,0.52) {$\vec E_0$};
  \draw[->,semithick,black!65] (0.40,-0.04)--(-0.40,-0.04);
  \node at (0,-0.24) {$\vec E_{\mathrm{ind}}$};
  \node at (0,-0.59) {$\vec E_{\mathrm{int}}=\vec0$};

  \draw[->,black!55] (-3.65,1.20)--(-3.12,1.20) node[right] {$x$};
  \draw[->,black!55] (-3.65,1.20)--(-3.65,1.73) node[above] {$y$};
  \node[align=center] at (0,-2.48)
    {Curvas exteriores: campo total\\[3pt]
     $\vec E_{\mathrm{tot}}=\vec E_0+\vec E_{\mathrm{ind}}$};
\end{tikzpicture}
\caption{Corte transversal del cilindro de radio $R$. Las flechas
superiores muestran el campo externo uniforme; las curvas exteriores
son líneas del campo total. La banda superficial representa la carga
inducida, negativa a la izquierda y positiva a la derecha. Las flechas
interiores muestran las contribuciones por separado:
$\vec E_{\mathrm{ind}}=-\vec E_0$, por lo que su suma es cero.}
\label{fig:cilindro-polarizado}
\end{figure}

\subsection*{Solución Detallada}

\subsubsection*{1. Planteamiento y solución general de Laplace}
En equilibrio, $\vec{E}_{\mathrm{int}}=\vec{0}$ y el cilindro es
equipotencial. Elegimos ese potencial como referencia:
\[
 V_{\mathrm{int}}=0,\qquad V_{\mathrm{ext}}(R,\phi)=0.
\]
En el vacío exterior, $s>R$, Poisson se reduce a Laplace:
\[
 \nabla^2V=-\frac{\rho}{\varepsilon_0},\qquad \rho=0
 \quad\Longrightarrow\quad \nabla^2V=0.
\]
Como no hay dependencia en $z$, usamos la forma de la solución de
Laplace estudiada en la ayudantía 5:
\begin{equation}
 \begin{aligned}
 V_{\mathrm{ext}}(s,\phi)
 &=A_0+B_0\ln s\\
 &\quad+\sum_{m=1}^{\infty}
 (A_ms^m+B_ms^{-m})
 [C_m\cos(m\phi)+D_m\sin(m\phi)].
 \end{aligned}
 \label{eq:cilindro-general}
\end{equation}
% La longitud de referencia del logaritmo queda absorbida en A_0.
En la suma, $m\in\mathbb{N}=\{1,2,3,\ldots\}$ es un número natural
que identifica el modo angular.

\subsubsection*{2. Determinación de Coeficientes}
El campo externo es $\vec{E}_0=E_0\hat{x}$. Un potencial que lo genera es
\[
 V_0=-E_0x,
 \qquad -\frac{\partial V_0}{\partial x}=E_0.
\]
Como $x=s\cos\phi$, lejos del cilindro debemos reproducir
\begin{equation}
 V_0(s,\phi)=-E_0s\cos\phi.
 \label{eq:potencial-uniforme}
\end{equation}
Comparamos \eqref{eq:potencial-uniforme} con
\eqref{eq:cilindro-general}: para obtener $s\cos\phi$ necesitamos
\[
 m=1,\qquad A_1C_1=-E_0.
\]
Los términos $s^m$ con $m>1$ crecerían más rápido que el dato impuesto;
por simetría no aparecen senos, y la neutralidad elimina el logaritmo:
\[
 A_m=0\ \ (m>1),\qquad D_m=0\ \ (m\geq1),\qquad B_0=0.
\]
La condición $V_{\mathrm{ext}}(R,\phi)=0$ descarta también los modos
decrecientes con $m>1$ y fija $A_0=0$. Absorbemos $C_1$ en la amplitud
radial, de modo que $B_1$ denota el coeficiente de $s^{-1}\cos\phi$:
\[
 V_{\mathrm{ext}}(s,\phi)
 =-E_0s\cos\phi+\frac{B_1}{s}\cos\phi.
\]
Evaluamos en la superficie:
\[
 -E_0R\cos\phi+\frac{B_1}{R}\cos\phi=0
 \quad\Longrightarrow\quad B_1=E_0R^2.
\]
Por tanto,
\begin{equation}
 \boxed{V_{\mathrm{ext}}(s,\phi)
 =-E_0\left(s-\frac{R^2}{s}\right)\cos\phi,\qquad s>R.}
 \label{eq:cilindro-potencial}
\end{equation}
\textbf{Las condiciones de frontera seleccionan los modos de la
solución general.}

\subsubsection*{3. Campo Eléctrico y Carga Inducida}
El campo total es $\vec E=\vec E_0+\vec E_{\mathrm{ind}}$, donde
$\vec E_{\mathrm{ind}}$ es el campo generado por la carga superficial.
Dentro del conductor, $\vec E_{\mathrm{ind}}=-\vec E_0$ y la suma es cero,
como muestra la figura~\ref{fig:cilindro-polarizado}.

En coordenadas cilíndricas y sin dependencia en $z$,
\[
 \vec{E}=-\nabla V
 =-\frac{\partial V}{\partial s}\hat{s}
  -\frac1s\frac{\partial V}{\partial\phi}\hat{\phi}.
\]
Al derivar \eqref{eq:cilindro-potencial}, obtenemos
\begin{equation}
 \boxed{E_s=E_0\left(1+\frac{R^2}{s^2}\right)\cos\phi,
 \qquad E_\phi=-E_0\left(1-\frac{R^2}{s^2}\right)\sin\phi.}
 \label{eq:cilindro-campo}
\end{equation}
En la superficie,
\[
 E_s(R,\phi)=2E_0\cos\phi,\qquad E_\phi(R,\phi)=0.
\]
El campo es normal a la superficie. Como $\vec{E}_{\mathrm{int}}=\vec{0}$,
la densidad inducida es
\begin{equation}
 \boxed{\sigma(\phi)=\varepsilon_0E_s(R,\phi)
 =2\varepsilon_0E_0\cos\phi.}
 \label{eq:cilindro-sigma}
\end{equation}
El lado derecho acumula carga positiva y el izquierdo, negativa.
Su distribución anula el campo externo dentro del conductor.

\clearpage
\section*{Visualización numérica}
\subsection*{1. Carga frente al plano conductor}
Las figuras evalúan las expresiones analíticas del ejercicio 1; no se
resuelve una ecuación diferencial de manera aproximada. Se usan las
coordenadas $x/d$, $y/d$ y $z/d$, junto con las cantidades adimensionales
\[
 \widetilde V=\frac{4\pi\varepsilon_0d}{q}V,\qquad
 \widetilde E=\frac{4\pi\varepsilon_0d^2}{q}|\vec{E}|,\qquad
 \widetilde\sigma=\frac{2\pi d^2}{q}\sigma.
\]
En los mapas de campo y potencial solo se representa $z>0$. Un disco
pequeño alrededor de la carga se excluye para evitar la singularidad del
modelo puntual. El error relativo del último panel compara el potencial
exacto con $V_{\mathrm{dip}}=2qd\cos\theta/(4\pi\varepsilon_0r^2)$.

\begin{figure}[!htbp]
\centering
\graficanumerica[width=0.98\textwidth]{carga_plano_visualizacion}
\caption{(a) Potencial exacto en el corte $y=0$, con $V=0$ en el
conductor. (b) Líneas del campo eléctrico; el color indica su magnitud
con escala logarítmica. (c) Densidad superficial inducida en $z=0$,
negativa y más intensa bajo la carga. (d) Error de la aproximación
dipolar para distintos ángulos: disminuye al aumentar $r/d$. Los mapas
usan expresiones exactas; el último panel compara el potencial exacto
con el término dipolar dominante.}
\label{fig:numerica-plano}
\end{figure}

\clearpage
\subsection*{2. Cilindro conductor en un campo uniforme}
Se evalúan las expresiones del ejercicio 2 en las coordenadas $x/R$,
$y/R$ y $z/R$. El potencial y el campo se normalizan por $E_0R$ y $E_0$,
respectivamente; la densidad se representa mediante
\[
 \widetilde\sigma=\frac{\sigma}{2\varepsilon_0E_0}=\cos\phi.
\]
El interior del conductor tiene $V=0$ y $\vec{E}=\vec{0}$. La representación
tridimensional muestra un tramo de un cilindro infinito: la densidad no
depende de $z$ y no se han añadido cargas en tapas ficticias.

\begin{figure}[!htbp]
\centering
\graficanumerica[width=0.98\textwidth]{cilindro_uniforme_visualizacion}
\caption{(a) Potencial en un corte transversal. (b) Líneas del campo
total y su magnitud: son normales a la superficie y se acercan a un
campo uniforme lejos del cilindro. (c) Densidad inducida sobre la
superficie cilíndrica; la escala de color distingue las zonas positivas
y negativas. (d) Perfil angular de la misma densidad. Los dos lados
acumulan cargas opuestas y la carga neta por unidad de longitud es cero.}
\label{fig:numerica-cilindro}
\end{figure}

\end{document}
